\documentclass[10pt,a4paper]{amsart}
\usepackage[margin=19mm]{geometry}
\usepackage{amsmath,amssymb,mathtools}
\usepackage[scr=rsfs,cal=euler]{mathalfa}
\usepackage{xfrac}
\usepackage[unicode,hidelinks,bookmarks=false]{hyperref}
\newcommand{\ie}{i.e.\ }
\newcommand{\cf}{cf.\ }
\newcommand{\resp}{respectively\ }
\newcommand{\Subsection}{\subsection}
\providecommand{\nobreakdash}{\nobreak\hbox{-}}
\setlength{\emergencystretch}{2em}
\multlinegap0pt
\usepackage{xcolor}
\usepackage{amssymb}
\newtheorem{proposition}{Proposition}
\title{On the membership of two-variable rational inner functions in spaces of Dirichlet type}
\author{Athanasios Beslikas, Alan Sola}
\date{Source: arXiv:2512.17388v2 (CC BY 4.0)}
\begin{document}
\maketitle

\noindent\emph{Source and licence.} On the membership of two-variable rational inner functions in spaces of Dirichlet type. Version arXiv:2512.17388v2 (CC BY 4.0), distributed by arXiv under the Creative Commons Attribution 4.0 International licence, http://creativecommons.org/licenses/by/4.0/, as recorded in the arXiv licence metadata for this exact version and shown on the abstract page https://arxiv.org/abs/2512.17388v2. Full licence notice: \texttt{licenses/arxiv-2512.17388-CC-BY-4.0.txt}.\par\smallskip

\noindent\textit{Editorial scope.} Counted unit: the article's proposition 5.1. (source0.tex lines 468--471). The article's complete original proof of it is delivered with it (source0.tex lines 472--485, one \texttt{proof} environment of 14 lines). No mathematics is written, repaired or re-derived: the statement and the proof are the article's own lines, in the article's own notation, and no retained line is substituted or re-flowed.  No further source-local context is needed: every cross-reference of every retained line resolves inside the two delivered spans.  Nothing was omitted from any retained span, so the package carries no omission record.  No retained line contains a citation, so the wrapper carries no bibliography and no bibliography entry.  No retained line contains a figure environment, an image-inclusion command or drawing code, so no image is delivered and the package declares no asset dependency.  Editorial formatting only, in addition: an AMS \texttt{amsart} preamble that restates the article's own declarations and macros used by the retained lines, namely the environments \texttt{proposition} on separate counters, together with the standard packages those lines need.  The retained environments print in this document as Proposition 1 in order of appearance, whereas the article prints the same items with the numbering of its own full document.  The complete native source of the article as distributed in the arXiv e-print for this exact version is bundled unchanged as \texttt{sources/191310/source0.tex}.
\begin{proposition} Let $\varphi=\frac{\tilde{p}}{p}$ a RIF and $L_j$ its refined Agler Kernels. Then $\varphi\in \mathcal{D}_{1,1,\alpha}$ if and only if
$|\zeta_j-\eta_j|^{\frac{1-\alpha}{2}}L_{j}(\zeta_1,\zeta_2,\eta_j)\in L^2(\mathbb{T}^3), j=1,2.$
\label{5.1.}
\end{proposition}

\begin{proof} The main tool for the proof of this result is the one dimensional weighted Douglas Formula. To see this, observe that 
$$\int_{\mathbb{T}}\int_{\mathbb{D}}|\partial_{z_1}\varphi(z_1,e^{i\theta})|^2(1-|z_1|^2)^{1-\alpha}dA(z_1)d\theta=\int_{\mathbb{T}}\|\varphi_{\zeta_2}(z_1)\|^2_{\mathcal{D}_{\alpha}(\mathbb{D})}d\theta.$$
It is well known that
$$\|\varphi_{\zeta_2}(z_1)\|^2_{\mathcal{D}_{\alpha}(\mathbb{D})}\approx\int_{\mathbb{T}^2}\frac{|\varphi_{\zeta_2}(\zeta_1)-\varphi_{\zeta_2}(\eta_1)|^2}{|\zeta_1-\eta_1|^{1+\alpha}}|d\zeta_1||d\eta_1|.$$
Now, we observe that by the refined Agler decomposition
$$\varphi_{\zeta_2}(\zeta_1)-\varphi_{\zeta_2}(\eta_1)=(\zeta_1-\eta_1)L_1(\zeta_1,\zeta_2,\eta_1),$$
$\zeta_1,\zeta_2,\eta_1 \in\mathbb{T}.$ Substituting into the one-dimensional Douglas Formula, we obtain
\begin{align}
\int_{\mathbb{T}}\int_{\mathbb{D}}|\partial_{z_1}\varphi(z_1,e^{i\theta})|^2(1-|z_1|^2)^{1-\alpha}dA(z_1)d\theta\approx &\int_{\mathbb{T}}\left(\int_{\mathbb{T}^2}\frac{|\varphi_{\zeta_2}(\zeta_1)-\varphi_{\zeta_2}(\eta_1)|^2}{|\zeta_1-\eta_1|^{1+\alpha}}|d\zeta_1||d\eta_1|\right)|d\zeta_2|&&\\
=&\int_{\mathbb{T}}\left(\int_{\mathbb{T}^2}\frac{|\zeta_1-\eta_1|^2|L_1(\zeta_1,\zeta_2,\eta_1)|^2}{|\zeta_1-\eta_1|^{1+\alpha}}\right)|d\zeta_2| \notag&&\\
=&\int_{\mathbb{T}^3}|\zeta_1-\eta_1|^{1-\alpha}|L_1(\zeta_1,\zeta_2,\eta_1)|^2|d\zeta_1||d\zeta_2||d\eta_1|\notag 
\end{align}
and the result follows.
\end{proof}

\end{document}
